\begin{exercises}
   \recommended \item 
     用高斯–若尔当消元求解每个方程组。
     \begin{exparts*}
        \partsitem \(
          \begin{linsys}[t]{2}
               x  &+  &y  &=  &2  \\
               x  &-  &y  &=  &0  
          \end{linsys}   \)
        \partsitem \(
          \begin{linsys}[t]{3}
               x  &   &   &-  &z  &=  &4  \\
              2x  &+  &2y &   &   &=  &1  
          \end{linsys}  \)
        \partsitem  \(
           \begin{linsys}[t]{2}
               3x  &-  &2y  &=  &1  \\
               6x  &+  &y   &=  &1/2 
           \end{linsys}  \)
        \partsitem \(
           \begin{linsys}[t]{3}
              2x  &-  &y  &  &  &= &-1  \\
               x  &+  &3y &- &z &= &5   \\
                  &   &y  &+ &2z&= &5   
           \end{linsys} \)
     \end{exparts*}
     \begin{answer}
       这些解答只给出了高斯–若尔当消元的过程。
       有了它，描述解集就很容易了。
       \begin{exparts}
         \partsitem 解集只含一个元素。
           \begin{multline*}
             \begin{amat}[r]{2}
               1  &1  &2  \\
               1  &-1 &0
             \end{amat}
             \grstep{-\rho_1+\rho_2}
             \begin{amat}[r]{2}
               1  &1  &2  \\
               0  &-2 &-2
             \end{amat}                          \\                       
             \grstep{-(1/2)\rho_2}
             \begin{amat}[r]{2}
               1  &1  &2  \\
               0  &1  &1
             \end{amat}                         
             \grstep{-\rho_2+\rho_1}
             \begin{amat}[r]{2}
               1  &0  &1  \\
               0  &1  &1
             \end{amat}
           \end{multline*}
         \partsitem 解集含有一个参数。
            \begin{equation*}
             \begin{amat}[r]{3}
               1  &0  &-1  &4  \\
               2  &2  &0   &1
             \end{amat}
             \grstep{-2\rho_1+\rho_2}
             \begin{amat}[r]{3}
               1  &0  &-1  &4  \\
               0  &2  &2   &-7
             \end{amat}                    
             \grstep{(1/2)\rho_2}
             \begin{amat}[r]{3}
               1  &0  &-1  &4  \\
               0  &1  &1   &-7/2
             \end{amat}
           \end{equation*}
         \partsitem 有唯一解。
           \begin{multline*}
             \begin{amat}[r]{2}
               3  &-2  &1  \\
               6  &1   &1/2
             \end{amat}
             \grstep{-2\rho_1+\rho_2}
             \begin{amat}[r]{2}
               3  &-2  &1  \\
               0  &5   &-3/2
             \end{amat}                            \\                       
             \grstep[(1/5)\rho_2]{(1/3)\rho_1}
             \begin{amat}[r]{2}
               1  &-2/3&1/3 \\
               0  &1   &-3/10
             \end{amat}                       
             \grstep{(2/3)\rho_2+\rho_1}
             \begin{amat}[r]{2}
               1  &0   &2/15 \\
               0  &1   &-3/10
             \end{amat}
          \end{multline*}
         \partsitem 第二步作一次对换两行的变换，会使算术更简单。
          \begin{multline*}
           \begin{amat}[r]{3}
             2  &-1  &0  &-1  \\
             1  &3   &-1 &5   \\
             0  &1   &2  &5
           \end{amat}
           \grstep{-(1/2)\rho_1+\rho_2}
           \begin{amat}[r]{3}
             2  &-1  &0  &-1   \\
             0  &7/2 &-1 &11/2 \\
             0  &1   &2  &5
           \end{amat}                                \\
           \begin{aligned}
           &\grstep{\rho_2\leftrightarrow\rho_3}
           \begin{amat}[r]{3}
             2  &-1  &0  &-1   \\
             0  &1   &2  &5    \\
             0  &7/2 &-1 &11/2
           \end{amat}                      
             \grstep{-(7/2)\rho_2+\rho_3}
             \begin{amat}[r]{3}
               2  &-1  &0  &-1   \\
               0  &1   &2  &5    \\
               0  &0   &-8 &-12
             \end{amat}                               \\     
             &\grstep[-(1/8)\rho_2]{(1/2)\rho_1}
             \begin{amat}[r]{3}
               1  &-1/2&0  &-1/2 \\
               0  &1   &2  &5    \\
               0  &0   &1  &3/2
             \end{amat}                                
             \grstep{-2\rho_3+\rho_2}
             \begin{amat}[r]{3}
               1  &-1/2&0  &-1/2 \\
               0  &1   &0  &2    \\
               0  &0   &1  &3/2
             \end{amat}                                    \\             
             &\grstep{(1/2)\rho_2+\rho_1}
             \begin{amat}[r]{3}
               1  &0   &0  &1/2  \\
               0  &1   &0  &2    \\
               0  &0   &1  &3/2
             \end{amat}
           \end{aligned}
         \end{multline*}
       \end{exparts}  
     \end{answer}
   \item 做高斯–若尔当消元。
     \begin{exparts*}
       \partsitem
         $\begin{linsys}[t]{3}
           x  &+ &y &- &z &= &3 \\
           2x &- &y  &- &z &= &1 \\
           3x &+  &y  &+ &2z &= &0
         \end{linsys}$
       \partsitem
         $\begin{linsys}[t]{3}
           x  &+ &y  &+ &2z  &= &0 \\
           2x  &- &y  &+  &z &= &1 \\
           4x &+ &y  &+ &5z &= &1  
         \end{linsys}$
     \end{exparts*}
     \begin{answer}
       \begin{exparts}
         \partsitem
           \begin{multline*}
             \begin{amat}{3}
               1 &1  &-1 &3 \\
               2 &-1 &-1 &1 \\
               3 &1  &2  &0
             \end{amat}
             \grstep[-3\rho_1+\rho_3]{-2\rho_1+\rho_2}
             \begin{amat}{3}
               1 &1  &-1 &3 \\
               0 &-3 &1  &-5 \\
               0 &-2  &5  &-9
             \end{amat}                             \\
           \begin{aligned}   
             &\grstep{-(2/3)\rho_2+\rho_3}
             \begin{amat}{3}
               1 &1  &-1    &3 \\
               0 &-3 &1     &-5 \\
               0 &0  &13/3  &-17/3
             \end{amat}                                 
             \grstep[(3/13)\rho_3]{-(1/3)\rho_2}
             \begin{amat}{3}
               1 &1  &-1    &3 \\
               0 &1  &-1/3    &5/3 \\
               0 &0  &1      &-17/13
             \end{amat}                                \\
             &\grstep[(1/3)\rho_3+\rho_2]{\rho_3+\rho_1}
             \begin{amat}{3}
               1 &1  &0    &22/13 \\
               0 &1  &0    &16/13 \\
               0 &0  &1      &-17/13
             \end{amat}
             \grstep{-\rho_2+\rho_1}
             \begin{amat}{3}
               1 &0  &0    &6/13 \\
               0 &1  &0    &16/13 \\
               0 &0  &1      &-17/13
             \end{amat}
           \end{aligned}
           \end{multline*}
         \partsitem
           \begin{multline*}
             \begin{amat}{3}
               1 &1  &2 &0 \\
               2 &-1 &1 &1 \\
               4 &1  &5 &1
             \end{amat}
             \grstep[-4\rho_1+\rho_3]{-2\rho_1+\rho_2}
             \begin{amat}{3}
               1 &1  &2  &0 \\
               0 &-3 &-3 &1 \\
               0 &-3 &-3 &1
             \end{amat}
             \grstep{-\rho_2+\rho_3}
             \begin{amat}{3}
               1 &1  &2  &0 \\
               0 &-3 &-3 &1 \\
               0 &0  &0  &0
             \end{amat}                              \\
             \grstep{-(1/3)\rho_2}
             \begin{amat}{3}
               1 &1  &2  &0 \\
               0 &1  &1 &-1/3 \\
               0 &0  &0  &0
             \end{amat}
             \grstep{-\rho_2+\rho_1}
             \begin{amat}{3}
               1 &0  &1  &1/3 \\
               0 &1  &1 &-1/3 \\
               0 &0  &0  &0
             \end{amat}
           \end{multline*}
       \end{exparts}
     \end{answer}
   \recommended \item 
     求出每个矩阵的简化阶梯形。
     \begin{exparts*}
       \partsitem \( \begin{mat}[r]
           2  &1  \\
           1  &3
         \end{mat}  \)
       \partsitem \( \begin{mat}[r]
           1  &3  &1  \\
           2  &0  &4  \\
          -1  &-3 &-3
         \end{mat}  \)
       \partsitem \( \begin{mat}[r]
           1  &0  &3  &1  &2  \\
           1  &4  &2  &1  &5  \\
           3  &4  &8  &1  &2
         \end{mat}  \)
       \partsitem \( \begin{mat}[r]
           0  &1  &3  &2  \\
           0  &0  &5  &6  \\
          1  &5  &1  &5
        \end{mat}  \)
    \end{exparts*}
    \begin{answer}
      使用高斯–若尔当消元。
      \begin{exparts}
        \partsitem 简化阶梯形除了由1构成的对角线以外全为0。
          \begin{equation*}
            \grstep{-(1/2)\rho_1+\rho_2}
            \begin{mat}[r]
              2  &1  \\
              0  &5/2
            \end{mat}
            \grstep[(2/5)\rho_2]{(1/2)\rho_1}
            \begin{mat}[r]
              1  &1/2\\
              0  &1
            \end{mat}                      
            \grstep{-(1/2)\rho_2+\rho_1}
            \begin{mat}[r]
              1  &0  \\
              0  &1
            \end{mat}
          \end{equation*}
        \partsitem 与上一题一样，简化阶梯形
          除了由1构成的对角线以外全是0。
          \begin{multline*}
            \grstep[\rho_1+\rho_3]{-2\rho_1+\rho_2}
            \begin{mat}[r]
              1  &3  &1  \\
              0  &-6 &2  \\
              0  &0  &-2
            \end{mat}
            \grstep[-(1/2)\rho_3]{-(1/6)\rho_2}
            \begin{mat}[r]
              1  &3  &1     \\
              0  &1  &-1/3  \\
              0  &0  &1
            \end{mat}                                      \\              
            \grstep[-\rho_3+\rho_1]{(1/3)\rho_3+\rho_2}
            \begin{mat}[r]
              1  &3  &0     \\
              0  &1  &0     \\
              0  &0  &1
            \end{mat}                  
            \grstep{-3\rho_2+\rho_1}
            \begin{mat}[r]
              1  &0  &0     \\
              0  &1  &0     \\
              0  &0  &1
            \end{mat}
           \end{multline*}
        \partsitem 列数比行数多，所以我们得到的
          不只是由1构成的一条对角线。
          \begin{multline*}
            \grstep[-3\rho_1+\rho_3]{-\rho_1+\rho_2}
            \begin{mat}[r]
              1  &0  &3  &1  &2  \\
              0  &4  &-1 &0  &3  \\
              0  &4  &-1 &-2 &-4
            \end{mat}
            \grstep{-\rho_2+\rho_3}
            \begin{mat}[r]
              1  &0  &3  &1  &2  \\
              0  &4  &-1 &0  &3  \\
              0  &0  &0  &-2 &-7
            \end{mat}                           \\
            \grstep[-(1/2)\rho_3]{(1/4)\rho_2}
            \begin{mat}[r]
              1  &0  &3    &1  &2  \\
              0  &1  &-1/4 &0  &3/4  \\
              0  &0  &0    &1  &7/2
            \end{mat}                           
            \grstep{-\rho_3+\rho_1}
            \begin{mat}[r]
              1  &0  &3    &0  &-3/2  \\
              0  &1  &-1/4 &0  &3/4     \\
              0  &0  &0    &1  &7/2
            \end{mat}
          \end{multline*}
        \partsitem 与上一小题一样，这不是一个方阵。 
          \begin{multline*}
            \grstep{\rho_1\leftrightarrow\rho_3}
            \begin{mat}[r]
              1  &5  &1  &5  \\
              0  &0  &5  &6  \\
              0  &1  &3  &2
            \end{mat}
            \grstep{\rho_2\leftrightarrow\rho_3}
            \begin{mat}[r]
              1  &5  &1  &5  \\
              0  &1  &3  &2  \\
              0  &0  &5  &6
            \end{mat}                    \\
            \begin{aligned}
              &\grstep{(1/5)\rho_3}
              \begin{mat}[r]
                1  &5  &1  &5  \\
                0  &1  &3  &2  \\
                0  &0  &1  &6/5
              \end{mat}                  
              \grstep[-\rho_3+\rho_1]{-3\rho_3+\rho_2}
              \begin{mat}[r]
                1  &5  &0  &19/5  \\
                0  &1  &0  &-8/5  \\
                0  &0  &1  &6/5
              \end{mat}                   \\
              &\grstep{-5\rho_2+\rho_1}
              \begin{mat}[r]
                1  &0  &0  &59/5  \\
                0  &1  &0  &-8/5  \\
                0  &0  &1  &6/5
              \end{mat}
            \end{aligned}
          \end{multline*}
      \end{exparts}  
    \end{answer}
  \item 求出每一个的简化阶梯形。
    \begin{exparts*}
      \partsitem $
        \begin{mat}
          0  &2 &1 \\
          2  &-1 &1 \\
          -2 &-1 &0
        \end{mat}$
      \partsitem $
        \begin{mat}
          1 &3 &1 \\
          2 &6 &2 \\
         -1 &0 &0
        \end{mat}$
    \end{exparts*}
    \begin{answer}
      \begin{exparts}
        \partsitem 先作对换。
          \begin{equation*}
            \grstep{\rho_1\leftrightarrow\rho_2}
            \grstep{\rho_1+\rho_3}
            \grstep{\rho_2+\rho_3}
            \grstep[(1/2)\rho_2 \\ (1/2)\rho_3]{(1/2)\rho_1}
            \grstep[(-1/2)\rho_3+\rho_2]{(-1/2)\rho_3+\rho_1}
            \grstep{(1/2)\rho_2+\rho_1}
            \begin{mat}
              1 &0 &0 \\
              0 &1 &0 \\
              0 &0 &1
            \end{mat}
          \end{equation*}
        \partsitem 这里对换在中间进行。
        \begin{equation*}
          \grstep[\rho_1+\rho_3]{-2\rho_1+\rho_2}
          \grstep{\rho_2\leftrightarrow\rho_3}
          \grstep{(1/3)\rho_2}
          \grstep{-3\rho_2+\rho_1}
          \begin{mat}
            1 &0 &0   \\
            0 &1 &1/3 \\
            0 &0 &0
          \end{mat}
        \end{equation*}
      \end{exparts}
    \end{answer}
  \recommended \item 
    用高斯–若尔当消元求出每个解集，
    然后读出参数化。
    \begin{exparts}
      \partsitem \( \begin{linsys}[t]{3}
                  2x  &+  &y  &-  &z  &=  &1  \\
                  4x  &-  &y  &   &   &=  &3  
                  \end{linsys}  \)
      \partsitem \( \begin{linsys}[t]{4}
                   x  &   &   &-  &z  &   &   &=  &1  \\
                      &   &y  &+  &2z &-  &w  &=  &3  \\
                   x  &+  &2y &+  &3z &-  &w  &=  &7  
                    \end{linsys}  \)
      \partsitem \( \begin{linsys}[t]{4}
                   x  &-  &y  &+  &z  &   &   &=  &0  \\
                      &   &y  &   &   &+  &w  &=  &0  \\
                  3x  &-  &2y &+  &3z &+  &w  &=  &0  \\
                      &   &-y &   &   &-  &w  &=  &0  
                  \end{linsys}  \)
      \partsitem \( \begin{linsys}[t]{5}
                   a  &+  &2b &+  &3c &+  &d  &-  &e  &=  &1  \\
                  3a  &-  &b  &+  &c  &+  &d  &+  &e  &=  &3  
                  \end{linsys}  \)
    \end{exparts}
    \begin{answer}
      至于高斯消元的那一半，
      参见第一章第~I.2 节习题
      \nearbyexercise{exer:SlvMatNot} 的解答。
      \begin{exparts}
      \partsitem “若尔当”那一半如下进行。
        \begin{equation*}
          \grstep[-(1/3)\rho_2]{(1/2)\rho_1}
          \begin{amat}[r]{3}
            1  &1/2 &-1/2 &1/2  \\
            0  &1   &-2/3 &-1/3
          \end{amat}                           
          \grstep{-(1/2)\rho_2+\rho_1}
          \begin{amat}[r]{3}
            1  &0   &-1/6 &2/3  \\
            0  &1   &-2/3 &-1/3
          \end{amat}
        \end{equation*}
        解集为下面这个集合
        \begin{equation*}
          \set{\colvec[r]{2/3 \\ -1/3 \\ 0}
               +\colvec[r]{1/6 \\ 2/3 \\ 1}z
              \suchthat z\in\Re}
        \end{equation*}
      \partsitem 后一半是
        \begin{equation*}
          \grstep{\rho_3+\rho_2}
          \begin{amat}[r]{4}
            1  &0  &-1  &0  &1 \\
            0  &1  &2   &0  &3 \\
            0  &0  &0   &1  &0
          \end{amat}
        \end{equation*}
        于是解如下。
        \begin{equation*}
          \set{\colvec[r]{1 \\ 3 \\ 0 \\ 0}
               +\colvec[r]{1 \\ -2 \\ 1 \\ 0}z
              \suchthat z\in\Re}
        \end{equation*}
      \partsitem 这个若尔当一半
        \begin{equation*}
          \grstep{\rho_2+\rho_1}
          \begin{amat}[r]{4}
            1  &0  &1   &1  &0 \\
            0  &1  &0   &1  &0 \\
            0  &0  &0   &0  &0 \\
            0  &0  &0   &0  &0
          \end{amat}
        \end{equation*}
        给出
        \begin{equation*}
          \set{\colvec[r]{0 \\ 0 \\ 0 \\ 0}
               +\colvec[r]{-1 \\ 0 \\ 1 \\ 0}z
               +\colvec[r]{-1 \\ -1 \\ 0 \\ 1}w
              \suchthat z,w\in\Re}
        \end{equation*}
        （当然，我们可以在描述中省去零向量）。
      \partsitem “若尔当”那一半
        \begin{align*}
          &\grstep{-(1/7)\rho_2}
          \begin{amat}[r]{5}
            1  &2  &3   &1   &-1   &1  \\
            0  &1  &8/7 &2/7 &-4/7 &0
          \end{amat}                   \\
          &\grstep{-2\rho_2+\rho_1}
          \begin{amat}[r]{5}
            1  &0  &5/7 &3/7 &1/7  &1  \\
            0  &1  &8/7 &2/7 &-4/7 &0
          \end{amat}
        \end{align*}
        以这个解集收尾。
        \begin{equation*}
          \set{\colvec[r]{1 \\ 0 \\ 0 \\ 0 \\ 0}
               +\colvec[r]{-5/7 \\ -8/7 \\ 1 \\ 0 \\ 0}c
               +\colvec[r]{-3/7 \\ -2/7 \\ 0 \\ 1 \\ 0}d
               +\colvec[r]{-1/7 \\ 4/7 \\ 0 \\ 0 \\ 1}e
              \suchthat c,d,e\in\Re}
        \end{equation*}
    \end{exparts}
   \end{answer}
  \item 
    给出这个矩阵的两个不同的阶梯形版本。
    \begin{equation*}
      \begin{mat}[r]
        2  &1  &1  &3  \\
        6  &4  &1  &2  \\
        1  &5  &1  &5
      \end{mat}
    \end{equation*}
    \begin{answer}
      按常规作高斯消元法便得到一个：
      \begin{equation*}
        \grstep[-(1/2)\rho_1+\rho_3]{-3\rho_1+\rho_2}
        \begin{mat}[r]
          2  &1  &1  &3  \\
          0  &1  &-2 &-7 \\
          0  &9/2&1/2&7/2
        \end{mat}
        \grstep{-(9/2)\rho_2+\rho_3}
        \begin{mat}[r]
          2  &1  &1    &3  \\
          0  &1  &-2   &-7 \\
          0  &0  &19/2 &35
        \end{mat}
      \end{equation*}
      而任何外观上的改动，比如把最下面一行乘以 \( 2 \)，
      \begin{equation*}
        \begin{mat}[r]
          2  &1  &1    &3  \\
          0  &1  &-2   &-7 \\
          0  &0  &19   &70
        \end{mat}
      \end{equation*}
      便给出另一个。  
    \end{answer}
  \recommended \item \label{exer:PossRedEchFrms} 
    列出各种尺寸可能出现的简化阶梯形。
    \begin{exparts*}
      \partsitem \( \nbyn{2} \)
      \partsitem \( \nbym{2}{3} \)
      \partsitem \( \nbym{3}{2} \)
      \partsitem \( \nbyn{3} \)
    \end{exparts*}
    \begin{answer}
      在下面列出的各情形中，我们取 $a,b\in\Re$。
      因此，下面列出的某些标准形
      实际上包含了无穷多种情形。
      特别地，它们包括了 $a=0$ 与 $b=0$ 的情形。
      \begin{exparts}
        \partsitem 
          $\begin{mat}[r]
            0  &0  \\
            0  &0
          \end{mat}$,
          $\begin{mat}[r]
            1  &a  \\
            0  &0
          \end{mat}$, 
          $\begin{mat}[r]
            0  &1  \\
            0  &0
          \end{mat}$, 
          $\begin{mat}[r]
            1  &0  \\
            0  &1
          \end{mat}$
        \partsitem
          $\begin{mat}[r]
               0  &0  &0  \\
               0  &0  &0
             \end{mat}$,
          $\begin{mat}[r]
               1  &a  &b  \\
               0  &0  &0
             \end{mat}$,
          $\begin{mat}[r]
               0  &1  &a  \\
               0  &0  &0
             \end{mat}$,
          $\begin{mat}[r]
               0  &0  &1  \\
               0  &0  &0
             \end{mat}$,
          $\begin{mat}[r]
               1  &0  &a  \\
               0  &1  &b
             \end{mat}$,
          $\begin{mat}[r]
               1  &a  &0  \\
               0  &0  &1
             \end{mat}$,
          $\begin{mat}[r]
               0  &1  &0  \\
               0  &0  &1
             \end{mat}$
        \partsitem
          $\begin{mat}[r]
               0  &0  \\
               0  &0  \\
               0  &0
             \end{mat}$,
          $\begin{mat}[r]
               1  &a  \\
               0  &0  \\
               0  &0
             \end{mat}$,
          $\begin{mat}[r]
               0  &1  \\
               0  &0  \\
               0  &0
             \end{mat}$,
          $\begin{mat}[r]
               1  &0  \\
               0  &1  \\
               0  &0
             \end{mat}$
        \partsitem
          $\begin{mat}[r]
               0  &0  &0  \\
               0  &0  &0  \\
               0  &0  &0
             \end{mat}$,
          $\begin{mat}[r]
               1  &a  &b  \\
               0  &0  &0  \\
               0  &0  &0
             \end{mat}$,
          $\begin{mat}[r]
               0  &1  &a  \\
               0  &0  &0  \\
               0  &0  &0
             \end{mat}$,
          $\begin{mat}[r]
               0  &1  &0  \\
               0  &0  &1  \\
               0  &0  &0
             \end{mat}$,
          $\begin{mat}[r]
               0  &0  &1  \\
               0  &0  &0  \\
               0  &0  &0
             \end{mat}$,
          $\begin{mat}[r]
               1  &0  &a  \\
               0  &1  &b  \\
               0  &0  &0
             \end{mat}$,
          $\begin{mat}[r]
               1  &a  &0  \\
               0  &0  &1  \\
               0  &0  &0
             \end{mat}$,
          $\begin{mat}[r]
               1  &0  &0  \\
               0  &1  &0  \\
               0  &0  &1
             \end{mat}$
      \end{exparts}  
    \end{answer}
  \recommended \item  
    对非奇异矩阵作高斯–若尔当消元
    会得到什么结果？
    \begin{answer}
      非奇异的齐次线性方程组有唯一解。
      所以非奇异矩阵必定能化简成一个（方） 
      矩阵，它除了从左上到右下的对角线上是 \( 1 \) 之外
      全为 \( 0 \)，
      比如下面这些。
      \begin{equation*}
         \begin{mat}[r]
           1  &0  \\
           0  &1  \\
         \end{mat}
         \qquad
         \begin{mat}[r]
           1  &0  &0  \\
           0  &1  &0  \\
           0  &0  &1
         \end{mat}
      \end{equation*}  
    \end{answer}
 \item 判断下述每个关系是否是 $\nbyn{2}$ 矩阵
   集合上的等价关系。
   \begin{exparts}
     \partsitem 若两个矩阵第一行第一列的元素相同，
       则称这两个矩阵相关。
     % \partsitem two matrices are related if their four entries sum to the 
     %   same total
     \partsitem 若两个矩阵第一行第一列的元素相同，
       或者第二行第二列的元素相同，
       则称这两个矩阵相关。
   \end{exparts}
   \begin{answer}
     \begin{exparts}
       \partsitem  这是一个等价关系。
       
          若两个矩阵相关，我们可以写成 $M_1\sim M_2$。
          $\sim$~关系是自反的，因为任何矩阵与它自身
          有相同的 $1,1$ 元素。
          该关系是对称的，因为如果 $M_1$ 与~$M_2$ 有相同的 $1,1$
          元素，那么显然 $M_2$ 与~$M_1$ 也有相同的
          $1,1$ 元素。
          最后，该关系是传递的，因为如果 $M_1\sim M_2$，即两者
          有相同的 $1,1$ 元素，
          且 $M_2\sim M_3$，即两者
          也有相同的元素，那么三者都有相同的~$1,1$ 元素，
          从而 $M_1\sim M_3$。 
       % \partsitem  This is an equivalence.
       %    Write $M_1\sim M_2$ if they are related.

       %    This relation is reflexive because any matrix has the 
       %    same sum of entries as itself.
       %    The $\sim$~relation is symmetric because if $M_1$ has the same
       %    sum of 
       %    entries as~$M_2$, then also $M_2$ has the same sum of entries
       %    as~$M_1$.
       %    Finally, the relation is transitive because if $M_1\sim M_2$ so their
       %    entry sum is the same, 
       %    and $M_2\sim M_3$ so they have the 
       %    same entry sum as each other, then all three have the same entry sum, 
       %    and $M_1\sim M_3$. 
       \partsitem
          这不是等价关系，因为它不是传递的。
          下面第一个矩阵与第二个矩阵因 $1,1$ 元素而相关，
          第二个与第三个因 $2,2$~元素而相关。
          但第一个与第三个不相关。
          \begin{equation*}
            \begin{mat}
              1 &0 \\
              0 &0
            \end{mat}
            \quad
            \begin{mat}
              1 &0 \\
              0 &-1
            \end{mat}
            \quad
            \begin{mat}
              0 &0 \\
              0 &-1
            \end{mat}
          \end{equation*}
     \end{exparts}
   \end{answer}
 \item \cite{Cleary}
    考虑 $\nbyn{2}$ 矩阵集合上的如下关系：  
    如果 $A$ 中所有元素之和与 $B$ 中所有元素之和相同，
    我们就说 $A$ 与 $B$ \textit{总和相像}。  
    例如，零矩阵与第一行有两个七、
    第二行有两个负七的矩阵
    总和相像。
    证明或否定这是 $\nbyn{2}$ 矩阵集合
    上的等价关系。
    \begin{answer}
      它是等价关系。
      为证明这一点，我们必须验证该关系
      是自反的、对称的和传递的。

      设所有矩阵都是 $\nbyn{2}$ 的。
      自反性：我们注意到矩阵与它自身
      的元素之和相同。
      对称性：设 $A$ 与~$B$ 的元素之和相同，
      那么显然 $B$ 与~$A$ 的元素之和也相同。
      传递性同样不难\Dash 如果 $A$ 与 $B$ 的元素之和相同，
      且 $B$ 与 $C$ 的元素之和相同，
      那么 $A$ 与 $C$ 的元素之和相同。
    \end{answer}
 \item \label{exer:INotJMakesRowOpsRev}
   \nearbylemma{le:RowOpsRev} 的证明中提到了
   倍加操作上的 $i\neq j$ 条件。
   \begin{exparts}
     \partsitem 写出一个元素都非零的 $\nbyn{2}$ 矩阵，
        并证明 $-1\cdot\rho_1+\rho_1$ 操作
        不会被 $1\cdot\rho_1+\rho_1$ 抵消。
     \partsitem 扩充该引理的证明，明确指出
        它究竟在何处用到倍加的 $i\neq j$ 条件。
   \end{exparts}
   \begin{answer}
    \begin{exparts}
      \partsitem 例如，
        \begin{equation*}
          \begin{mat}[r]
            1  &2  \\
            3  &4  
          \end{mat}
          \grstep{-\rho_1+\rho_1}
          \begin{mat}[r]
            0  &0  \\
            3  &4  
          \end{mat}
          \grstep{\rho_1+\rho_1}
          \begin{mat}[r]
            0  &0  \\
            3  &4  
          \end{mat}
        \end{equation*}
        矩阵被改变了。
      \partsitem 这个操作
        \begin{equation*}
          \begin{mat}
            \vdotswithin{a_{i,1}}                     \\
            a_{i,1}  &\cdots  &a_{i,n}  \\
            \vdotswithin{a_{i,1}}                     \\
            a_{j,1}  &\cdots  &a_{j,n}  \\
            \vdotswithin{a_{i,1}}                     
          \end{mat}
          \grstep{k\rho_i+\rho_j}
          \begin{mat}
            \vdotswithin{a_{i,1}}                                      \\
            a_{i,1}           &\cdots  &a_{i,n}          \\
            \vdotswithin{a_{i,1}}                                      \\
            ka_{i,1}+a_{j,1}  &\cdots  &ka_{i,n}+a_{j,n}  \\
            \vdotswithin{a_{i,1}}                     
          \end{mat}
        \end{equation*}      
        由于 $i\neq j$ 的限制而使第$i$行保持不变。
        因为第$i$行未变，所以这个操作
        \begin{equation*}                                        
          \grstep{-k\rho_i+\rho_j}
          \begin{mat}
            \vdotswithin{a_{i,1}}                                      \\
            a_{i,1}           &\cdots  &a_{i,n}          \\
            \vdotswithin{a_{i,1}}                                      \\
            -ka_{i,1}+ka_{i,1}+a_{j,1}  &\cdots &-ka_{i,n}+ka_{i,n}+a_{j,n} \\
            \vdotswithin{a_{i,1}}                     
          \end{mat}
        \end{equation*}
        把第$j$行还原到它原来的状态。
        % (If $i=j$ then the third matrix would have entries of the 
        % form $-k(ka_{i,j}+a_{i,j})+ka_{i,j}+a_{i,j}$.)
    \end{exparts}
   \end{answer}
 \recommended \item \cite{Cleary}
  考虑某个班级学生的集合。  
  下列关系中哪些是等价关系？  
  对每个答案至少用一句话解释。
  \begin{exparts}
    \item  两个学生 $x, y$ 相关，
      如果 $x$ 修过的数学课
      至少与 $y$ 一样多。
    \item 学生 $x, y$ 相关，如果他们的名字
      以同一个字母开头。
  \end{exparts}
  \begin{answer}
    要成为等价关系，每个关系都必须是自反的、对称的和
    传递的。
    \begin{exparts}
      \item 这个关系
        不是对称的，因为如果 $x$ 修了 $4$~门课而 $y$
        修了 $3$ 门，那么 $x$ 与 $y$ 相关，但 $y$ 与 $x$
        不相关。
      \item 这是自反的，因为 $x$ 的名字与 $x$ 自己的名字
        以同一个字母开头。
        它是对称的，因为如果 $x$ 的名字与 $y$ 的名字
        以同一个字母开头，那么 $y$ 的名字也与~$x$ 的以同一个字母开头。
        它还是传递的，因为如果 $x$ 的名字与~$y$ 的以同一个字母开头，
        且 $y$ 的名字与 $z$ 的
        以同一个字母开头，那么 $x$ 的名字与 $z$ 的以同一个字母开头。
        所以它是等价关系。
    \end{exparts}
  \end{answer}
  \item
  证明下述每一个都是 $\nbyn{2}$ 矩阵集合
  上的等价关系。
  描述各个等价类。
  \begin{exparts}
    \item 如果两个矩阵沿对角线的乘积相同，
      也就是说，左上角元素与
      右下角元素的乘积相等，则这两个矩阵相关。
    \item 如果两个矩阵都至少有一个元素
      为~$1$，
      或者两者都没有，则这两个矩阵相关。
  \end{exparts}
  \begin{answer}
    对每一个我们都必须验证自反性、对称性和
    传递性这三个条件。
    \begin{exparts}
      \item 显然，任何矩阵沿对角线的乘积都与它自身相同，
        所以该关系是自反的。
        该关系是对称的，因为如果 $A$ 沿对角线的乘积
        与~$B$ 的相同，
        即 $a_{1,1}\cdot a_{2,2}=b_{1,1}\cdot b_{2,2}$，
        那么 $B$ 的乘积也与~$A$ 的相同。
 
        传递性类似：设 $A$ 的乘积为~$r$，
        且等于 $B$ 的乘积。
        又设 $B$ 的乘积等于 $C$ 的乘积。
        那么三者的乘积都是~$r$，即 $A$ 的乘积等于~$C$ 的乘积。

        每个实数对应一个等价类，也就是说，
        该等价类包含所有沿对角线
        乘积为这个实数的 $\nbyn{2}$ 矩阵。
      \item 自反性：如果矩阵 $A$ 有一个元素为~$1$，那么它与
        自身相关；如果没有，那么它同样与
        自身相关。
        对称性也分两种情形：设 $A$ 与~$B$ 相关。
        如果 $A$ 有一个元素为~$1$，那么~$B$ 也有，于是 $B$ 与~$A$ 相关。
        如果 $A$ 没有~$1$，那么 $B$ 也没有，同样 $B$ 与
        $A$ 相关。
 
        传递性：设 $A$ 与~$B$ 相关，$B$ 与~$C$ 相关。
        如果 $A$ 有一个元素为~$1$，那么~$B$ 也有；又因 $B$ 与~$C$ 相关，
        所以~$C$ 也有，
        从而 $A$ 与~$C$ 相关。
        同样，如果 $A$ 没有~$1$，那么~$B$ 也没有，因而
        ~$C$ 也没有，结论同样是 $A$ 与~$C$ 相关。 

        恰有两个等价类：一个包含任何
        至少有一个元素为~$1$ 的 $\nbyn{2}$ 矩阵，
        另一个包含所有不含 $1$ 的矩阵。
    \end{exparts}
  \end{answer}
  \item 证明下述每一个都不是 $\nbyn{2}$ 矩阵
    集合上的等价关系。
    \begin{exparts}
      \item 两个矩阵 $A,B$ 相关，如果 $a_{1,1}=-b_{1,1}$。
      \item 如果两个矩阵的元素之和相差在 $5$ 以内，
        即当
        $\absval{(a_{1,1}+\cdots+a_{2,2})-(b_{1,1}+\cdots+b_{2,2})}<5$ 时 $A$ 与~$B$ 相关，则称这两个矩阵相关。
    \end{exparts}
    \begin{answer}
      \begin{exparts}
        \item 这个关系不是自反的。
          例如，任何左上角元素为~$1$ 的矩阵
          都与自身不相关。
        \item 这个关系不是传递的。
          对下面这三个矩阵，$A$ 与~$B$ 相关，$B$ 与~$C$ 相关，
          但 $A$ 与~$C$ 不相关。
          \begin{equation*}
            A=\begin{mat}
              0 &0 \\
              0 &0
            \end{mat},\quad
            B=\begin{mat}
              4 &0 \\
              0 &0
            \end{mat},\quad
            C=\begin{mat}
              8 &0 \\
              0 &0
            \end{mat},\quad
          \end{equation*}
      \end{exparts}
    \end{answer}
\end{exercises}




















